\chapter{Sai, quê và cảm giác thua rất thật}
\markboth{Sai, quê và cảm giác thua rất thật}{Sai, quê và cảm giác thua rất thật}

\section{Em ngại phát biểu vì sai là bị cười}
\begin{truyen}{Em ngại phát biểu vì sai là bị cười}{Classroom Talk}
\chuhoa{H}{ọc sinh: ``Em không phải không biết gì. Em chỉ sợ nói sai xong bị cười, rồi quê tới hết buổi.''}
\textit{Student: ``It is not that I know nothing. I am just afraid that if I say something wrong, people will laugh and I will feel embarrassed for the rest of the class.''}\\[4pt]

\teacherpair{Vậy em muốn quê một lần để biết đường sửa, hay dốt âm thầm lâu dài?}{So would you rather feel embarrassed once and know what to fix, or stay silently weak for a long time?}

\studentpair{Thầy nói câu nào cũng đâm trúng chỗ em né.}{Every sentence you say hits the exact spot I keep trying to avoid.}

\teacherpair{Vì chỗ em né lại chính là chỗ em cần bước qua.}{That is because the place you avoid is exactly the place you need to walk through.}
\end{truyen}

\begin{baihoc}
Nỗi quê ngắn hạn thường rẻ hơn cái giá của việc im lặng quá lâu.
\textit{(Short-term embarrassment often costs less than long-term silence.)}
\end{baihoc}

\begin{ghinhoanh}
\textbf{Short English to remember:}

Embarrassment passes.

Untouched weakness stays.
\end{ghinhoanh}

\begin{tuhocnhanh}
1. Em đang né học vì sợ sai hay sợ bị cười? \textit{(Are you avoiding learning because of error or embarrassment?)}

2. Em có thể tập phát biểu ở môi trường an toàn hơn trước không? \textit{(Can you practice speaking in a safer setting first?)}
\end{tuhocnhanh}
\ngancach

\section{Thua đứa mình từng giúp còn đau hơn}
\begin{truyen}{Thua đứa mình từng giúp còn đau hơn}{Classroom Talk}
\chuhoa{H}{ọc sinh: ``Thầy ơi, thua người giỏi sẵn em chịu. Nhưng thua đứa mình từng chỉ bài mới cay.''}
\textit{Student: ``Sir, I can accept losing to someone who was already good. But losing to someone I used to teach hurts much more.''}\\[4pt]

\teacherpair{Em cay vì nó giỏi lên, hay cay vì em tưởng nó phải đứng dưới mình mãi?}{Are you upset because that person improved, or because you thought they were supposed to stay below you forever?}

\studentpair{Chắc cả hai.}{Probably both.}

\teacherpair{Vậy em đau không chỉ vì kết quả. Em đau vì cái tôi bị trầy.}{Then you are not hurting only because of the result. You are hurting because your ego got scratched.}
\end{truyen}

\begin{baihoc}
Người khác tiến bộ không phải là phản bội mình. Đôi khi chỉ là mình đã đứng yên quá lâu.
\textit{(Another person improving is not a betrayal. Sometimes it only shows that we stood still too long.)}
\end{baihoc}

\begin{ghinhoanh}
\textbf{Short English to remember:}

Their growth is not your insult.

Pride often hurts before reality does.
\end{ghinhoanh}

\begin{tuhocnhanh}
1. Em đang buồn vì thua, hay vì cái tôi bị chạm? \textit{(Are you upset by the loss, or by wounded pride?)}

2. Em cần học gì từ người đã vượt lên? \textit{(What can you learn from the person who surpassed you?)}
\end{tuhocnhanh}
\ngancach

\section{Em sửa mấy lần vẫn sai nên muốn bỏ}
\begin{truyen}{Em sửa mấy lần vẫn sai nên muốn bỏ}{Classroom Talk}
\chuhoa{H}{ọc sinh: ``Em sửa đi sửa lại vẫn trật. Cảm giác như đầu em không hấp thụ nổi nữa. Muốn bỏ luôn cho xong.''}
\textit{Student: ``I keep correcting it and I am still wrong. It feels like my head cannot absorb any more. I just want to quit and be done with it.''}\\[4pt]

\teacherpair{Bỏ lúc này sẽ cho em cảm giác nhẹ ngay. Nhưng cái sai đó không tự biến mất. Nó chỉ chờ dịp khác quay lại đập em mạnh hơn.}{Quitting now will make you feel lighter right away. But that mistake will not disappear on its own. It will simply wait for another moment to hit you even harder.}

\studentpair{Vậy cứ lì tiếp hả thầy?}{So should I just stubbornly keep going?}

\teacherpair{Không phải lì mù. Là đổi cách sửa, đổi người hỏi, đổi góc nhìn. Đừng cứ dùng một kiểu rồi kết luận mình hết cứu.}{Not blindly. You need to change the method, ask a different person, or look at the problem from a different angle. Do not keep using one failed approach and then conclude that you are beyond help.}
\end{truyen}

\begin{baihoc}
Muốn bỏ thường đến ngay trước lúc cần đổi cách, không phải lúc cần tự kết thúc mình.
\textit{(The urge to quit often appears right when a change of method is needed, not a final surrender.)}
\end{baihoc}

\begin{ghinhoanh}
\textbf{Short English to remember:}

Repeating the same failed method is not resilience.

Change the method before judging yourself.
\end{ghinhoanh}

\begin{tuhocnhanh}
1. Em đã thử đổi cách học ở chỗ sai đó chưa? \textit{(Have you changed your method for that weak point?)}

2. Ai có thể giúp em nhìn ra lỗi theo cách khác? \textit{(Who can help you see the mistake differently?)}
\end{tuhocnhanh}
\ngancach

\section{Có lúc em thấy mình cố cũng chỉ vậy}
\begin{truyen}{Có lúc em thấy mình cố cũng chỉ vậy}{Classroom Talk}
\chuhoa{H}{ọc sinh: ``Thầy ơi, có lúc em thấy mình cố rồi mà đời vẫn không đổi bao nhiêu. Nên em tự hỏi cố nữa để làm gì.''}
\textit{Student: ``Sir, sometimes I feel like I have already tried, but life barely changes. So I keep asking myself why I should continue.''}\\[4pt]

\teacherpair{Có những giai đoạn cố gắng chưa ra quả ngay. Nhưng nó đang âm thầm dời cái nền của em. Vấn đề là mình có đủ kiên nhẫn để chịu đoạn chưa thấy gì không.}{There are seasons when effort does not show results immediately. But that effort is still quietly shifting your foundation. The question is whether you have enough patience to endure the part where nothing seems visible yet.}

\studentpair{Đoạn đó khó chịu thật.}{That part really does feel uncomfortable.}

\teacherpair{Ừ. Trưởng thành là vẫn làm việc đúng trong lúc chưa được thưởng.}{Yes. Maturity is doing the right work even while no reward has arrived yet.}
\end{truyen}

\begin{baihoc}
Cái khó nhất không phải cố một bữa. Là cố qua giai đoạn chưa được công nhận.
\textit{(The hardest part is not trying once. It is trying through the phase where no one notices.)}
\end{baihoc}

\begin{ghinhoanh}
\textbf{Short English to remember:}

Invisible effort still matters.

Growth often starts below the surface.
\end{ghinhoanh}

\begin{tuhocnhanh}
1. Em có đang bỏ một việc đúng chỉ vì chưa thấy kết quả sớm? \textit{(Are you dropping the right work just because results are slow?)}

2. Em cần thêm bằng chứng nào để tin vào quá trình của mình? \textit{(What evidence would help you trust your process more?)}
\end{tuhocnhanh}
\ngancach